\subsection{معرفی و انگیزه}
این ماژول در چارچوب پژوهش حاضر پیاده‌سازی و ارزیابی شده و فاقد مقاله مستقل جداگانه در ادبیات تحقیق است.

مکانیزم توجه \lr{Polynomial Attention with Dropout Regularization (PADRe)} به عنوان جایگزینی باثبات و کارآمد برای عملگر غیرخطی سنتی \lr{Softmax} در معماری ترنسفورمر طراحی شده است. تابع نمایی در توجه استاندارد اگرچه به تفکیک شدید توزیع توجه کمک می‌کند، اما از منظر محاسباتی هزینه‌بر بوده و محاسبه مخرج جمع نمایی به همراه محاسبات ممیز شناور، امکان فاکتورسازی هسته‌ای خطی را با دشواری روبه‌رو می‌سازد. علاوه بر این، در آموزش نیم‌دقت (\lr{FP16/BF16})، مقادیر نمایی ممکن است دچار سرریز (\lr{Overflow}) یا فروریزش (\lr{Underflow}) شوند.

ایده محوری در ماژول \lr{PADRe}، استفاده از یک کِرنل چندجمله‌ای متمرکز و کنترل‌شده به فرم $(1 + S/\text{order})^{\text{order}}$ با تنظیم‌گری دراپ‌اوت (\lr{Dropout Regularization}) است. این چندجمله‌ای تقریب بسیار دقیقی از رفتار نمایی در نزدیکی نقطه مبنا ارائه داده و هم‌زمان با مهار بیشینه‌ها از طریق مرکزسازی (\lr{Max-Centering})، امکان پایدارسازی عددی کامل، ادغام با سازوکار سرهای پرس‌وجوی گروه‌بندی‌شده (\lr{GQA}) و بهره‌گیری از آداپتورهای کم‌مرتبه تطبیقی (\lr{LoRA Residual Adapters}) را بدون برهم‌زدن وزن‌های پایه مدل از پیش‌آموزش‌دیده فراهم می‌آورد.
